%%-*-latex-*-

\begin{frame}
\frametitle{Basic definitions}

An \textbf{alphabet} is a non\hyp{}empty, finite set of symbols, called
\textbf{letters}, for example \word{a}, \word{b} etc.

\bigskip

A \textbf{word} is a finite series of letters from an alphabet, like
\word{abracadabra}. A letter is a word.

\bigskip

Just as letters can be joined to make up words, so can words: the word
\(u \cdot v\) is made of the letters of word~\(u\) followed by the
letters of word~\(v\), for instance, if \(u=\word{back}\) and
\(v=\word{up}\), then \(u \cdot v = \word{backup}\). This operation is
called \textbf{concatenation}.

\bigskip

Concatenation is associative: \((u \cdot v) \cdot w = u \cdot (v \cdot
w)\). The operator may be omitted: \((uv)w = u(vw)\). Concatenation is
a non\hyp{}commutative product, so it has a neutral element
\(\varepsilon\), called the \textbf{empty word}: \(u \cdot \varepsilon
= \varepsilon \cdot u = u\).

\end{frame}


\begin{frame}
\frametitle{Basic definitions}

We denote the \textbf{repetition} of a word with an exponent, e.g.
\(\word{a}^3 = \word{aaa}\). The \textbf{length} of a string \(x\) is
the length of the associated sequence and is noted \(\len{x}\). So
\(\len{\word{a}^3} = 3\).

\bigskip

A word~\(x\) is a \textbf{factor} of a word~\(y\) if there exists two
words \(u\)~and~\(v\) (possibly empty) such that \(y = uxv\). In other
words, \(u\)~occurs in~\(y\).

\bigskip

The word~\(x\) is a \textbf{prefix} of~\(y\), noted \(x \prefeq y\),
if~\(u=\varepsilon\), that is, if \(y = xv\). Moreover, it is a
\textbf{proper prefix}, noted \(x \pref y\), if \(v \neq \varepsilon\).

\bigskip

Given \(y = uxv\), the word~\(x\) is a \textbf{suffix} of~\(y\) if \(v
= \varepsilon\). Furthermore, it is a \textbf{proper suffix} if \(u
\neq \varepsilon\).

\bigskip

Let \(a\)~range over letters, and \(x\) and \(y\) over words. The
prefix relation is easy to define inductively:
\begin{itemize}

  \item \(\forall y. \, \varepsilon \prefeq y\);

  \item \(\forall a, x, y. \, x \prefeq y \Rightarrow a \cdot x \prefeq a
    \cdot y\).

\end{itemize}

\end{frame}

\begin{frame}
\frametitle{Basic definitions}

A letter in a word can be uniquely characterised by a natural number,
called \textbf{index}, assuming that the first letter has
index~\(0\). Read
{\small\url{https://www.cs.utexas.edu/~EWD/transcriptions/EWD08xx/EWD831.html}}

\bigskip

For example, if \(x=\word{top}\), then the letter at index~\(0\) is
written \(\ind{x}{0}=\word{t}\) and the one at index~\(2\) is
\(\ind{x}{2}=\word{p}\).

\bigskip

A factor~\(x\) of~\(y\) can be identified by the index
of~\(\ind{x}{0}\) in~\(y\).

\bigskip

The end of the factor can also be given; for example, \(x=\word{sit}\)
is a factor of \(y=\word{curiosity}\) at index~\(5\), written
\(\ind{y}{5,7} = x\). This is sometimes called a \textbf{slice} in
some programming languages. Sometimes the upper bound can be left
unspecified: \(\ind{y}{i\ldots}\), meaning \(\ind{y}{i,\len{y}-1}\),
if \(i < \len{y}\).

\end{frame}


\begin{frame}
\frametitle{The prefix predicate}

An imperative definition:

\begin{codebox}
\(\proc{Prefix} (x, t)\)\\
\li \(i \gets 0\)
\li \While \(i < \len{x}\) \LogAnd \(i < \len{t}\)
\li \Do \If \(x[i] = t[i]\)
\li \Then \(i \gets i + 1\)
\li \Else \(i \gets \len{x} + 1\) \End \End
\li \(\proc{Prefix} \gets i = \len{x}\) \End
\end{codebox}

Does it terminate? What test cases would be meaningful? Is it correct?
Is it optimal?

\end{frame}


\begin{frame}
\frametitle{The prefix predicate}

A single static assignment definition:

\begin{codebox}
\(\proc{Prefix} (x, t)\)\\
\li \(\proc{Prefix} \gets \proc{Prefix'}(0,x,t)\)
\end{codebox}

\begin{codebox}
\(\proc{Prefix'} (i,x, t)\)\\
\li \If \(i < \len{x}\) \LogAnd \(i < \len{t}\)
\li \Then \If \(x[i] = t[i]\)
\li       \Then \(\proc{Prefix'}(i + 1, x, t)\)
\li       \Else \(\proc{Prefix'}(\len{x} + 1, x, t)\) \End
\li \Else \(\proc{Prefix'} \gets i = \len{x}\) \End
\end{codebox}

\end{frame}

\begin{frame}
\frametitle{Na\"{\i}ve search}

Factor matching is common in text editing, although it is usually
better known as \textbf{exact string matching} in the academic field
of \emph{text algorithmics} (sometimes called \emph{stringology}).

\bigskip

Because of the asymmetric nature of the search, the word~\(p\) to find
will from now on be called the \textbf{pattern}, and the word~\(t\)
the \textbf{text}.

\bigskip

Assuming \(p\)~and~\(t\) are not empty, the na\"{\i}ve search simply
consists in checking whether \(p\)~is a prefix of~\(t\): if so, we are
done, otherwise, the pattern is shifted by one letter to the right
and, again, it is checked whether it is a prefix of that text suffix.

\bigskip

If \(p\)~cannot be shifted anymore because its end would surpass the
end of~\(t\), then it is not a factor: the search failed.

\end{frame}


\begin{frame}
  \frametitle{Na\"{\i}ve search}

\begin{figure}
\centering
\includegraphics[bb=75 621 352 715]{naive}
\end{figure}
Na\"{\i}vely matching pattern~\(p\) against text~\(t\)
  (failure in grey), where \(\ind{p}{i} \neq \ind{t}{j}\) (the letters
`\word{a}'~and~`\word{b}' are not relevant in themselves).

\end{frame}


\begin{frame}
  \frametitle{Na\"{\i}ve search}

  An imperative program:

\begin{codebox}
\(\proc{Naive} (p, t)\)\\
\li \(i \gets 0\); \(j \gets 0\)
\li \While \(i < \len{p}\) \LogAnd \(j < \len{t}\)
\li \Do \If \(t[j] = p[i]\)
\li \Then \(j \gets j + 1\); \(i \gets i + 1\)
\li \Else \(j \gets j + 1 - i\); \(i \gets 0\)\>\>\>\>\>\>\Comment We
slide \(p\) by 1. \End \End
\li \(\proc{Naive} \gets i = \len{p}\)
\end{codebox}

What test cases would be meaningful?

\bigskip

Would it be clearer to use two loops?

\bigskip

How about a version with single static assignments and slices? Without
slices?

\end{frame}

\begin{frame}
\frametitle{Na\"{\i}ve search/Analysis}

The input is a pattern~\(p\) of length~\(m\) and a text~\(t\) of
length~\(n\)

\bigskip

What configuration of the input leads to the minimal number of
comparisons? What is that number?

\bigskip

In the worst case, \(p\) is not in \(t\) and the test at line \(3\)
always fails on the last letter of \(p\), leading to make \(p\)
slide of one position at line \(5\) the latest as possible.

\bigskip

An example of positive match in the worst case is \(p =
\texttt{a}^{m-1}\texttt{b}\) and \(t = \texttt{a}^{n-1}\texttt{b}\),
with \(m \leq n\).

\bigskip

There are \(n - m + 1\) positions in \(t\) to compare with the first
letter of \(p\), and since there are \(m\) letters in \(p\), it makes
a total of \((n - m + 1) m = nm - m^2 + m < \len{t} \times \len{p}\)
positions.

\bigskip

What about the worst case when there is no match?

\end{frame}

\begin{frame}
  \frametitle{Morris-Pratt algorithm}

In case of mismatch, the na\"{\i}ve algorithm resumes comparing the
first letters of~\(p\) without using the information of the partial
success, to wit, we know \(\ind{p}{0,i-1} = \ind{t}{j-i,j-1}\) and
\(\ind{p}{i} \neq \ind{t}{j}\).

\bigskip

If we know an index~\(k\) such that \(\ind{p}{0,k-1} =
\ind{p}{i-k,i-1}\), that is, \(\ind{p}{0,k-1}\)~is a \emph{border}
of~\(\ind{p}{0,i-1}\) (also known as a \emph{side}), then we can
resume by comparing~\(\ind{t}{j}\) with~\(\ind{p}{k}\).

\bigskip

Clearly, the greater~\(k\), the more comparisons are skipped, so we
want to find the \textbf{maximum borders} of the prefixes of~\(p\).

\end{frame}

\begin{frame}
  \frametitle{Morris-Pratt algorithm}

\begin{figure}
\centering
\includegraphics[bb=74 621 367 717]{mp}
\end{figure}

where \(\Border{}{p}\) is the maximum border of~\(p\).

\bigskip

Contrary to na\"{\i}ve factoring, letters in the text are compared in
a strictly increasing order (never having to backtrack).

\end{frame}

\begin{frame}
  \frametitle{Morris-Pratt algorithm}

The border of a non\hyp{}empty word~\(y\) is a proper prefix of~\(y\)
which is also a suffix.

\bigskip

For example, the word `\word{abacaba}' has three borders:
\(\varepsilon\), `\word{a}'~and~`\word{aba}'. The last one is the
maximum border and we write
\(\Border{}{\underline{\word{aba}}\word{c}\underline{\word{aba}}} =
\word{aba}\).

\bigskip

Another example is \(\Border{}{\word{abac}} = \varepsilon\), simply
because we have \(\word{abac} =
\underline{\varepsilon}\word{abac}\underline{\varepsilon}\).

\bigskip

Note that maximum borders can overlap; consider, for example,
\(\Border{}{\underline{\word{aaa}}\word{a}} =
\Border{}{\word{a}\underline{\word{aaa}}} = \word{aaa}\).

\end{frame}

\begin{frame}
  \frametitle{Morris-Pratt algorithm}

Let us consider a run ending in a failure:

\begin{figure}
\centering
\includegraphics{mp_ex}
\end{figure}

\end{frame}
